\scanpage{9}
Si $E$ est un \deleted{motif}\inserted{$\unsure{Q|\mathbb{Q}}_{\ell}$-faisceau} sur $\unsure{\operatorname{Spec}|\operatorname{Sp}}\unsure{(\mathbb{F}_p,)|(\mathbb{F}_p),}$ \deleted{\unsure{sur|dans}} on trouve donc
\[
L_{\check E}(t)=(-t)^{-n}\det f_E\ L_E(t^{-1})
\]
où
\[
\left\{\begin{aligned}
n&=\operatorname{rang} E
  &&\text{et s'il s'agit d'un \uline{motif} \inserted{de poids $\rho$}}\\
 &&&\text{satisfaisant aux conj. de Weil}\\
\det f_E&=\deleted{\unsure{\varepsilon}}\inserted{\varepsilon(E)}\deleted{\unsure{\pm|+}}p^{\frac{n\rho}{2}}
\end{aligned}\right.
\]
avec $\varepsilon=\pm1$ [dépendant du motif $\bigwedge^{\unsure{n|{-}}}E$, \unsure{savoir} si $f$ y opère \unsure{positivement|parfaitement} ou non] ; noter que $\frac{n\rho}{2}$ est tjrs entier (i.e. $\rho$ impair $\Longrightarrow$ $n$ pair).

On trouve donc
\[
L_{\check E}(t)=\varepsilon(E)(-p^{\frac{\rho}{2}}/t)^n L_E(t^{-1})\ .
\]

D'autre part, par Lefschetz, on aura
\[
\check E\simeq E(\rho)
\]
d'où
\[
\begin{aligned}
L_{\check E}(t)&=\bigl(L_E(t)\ast L_{\mathbb{Q}_{\ell}(\rho)}(t)\bigr)^{-1}\\
&=\left(L_E(t)\ast\frac{1}{1-p^{-\rho}t}\right)^{-1}=L_E(t)\ast(1-p^{-\rho}t)
\end{aligned}
\]
\[
L_{\check E}(t)=L_E\left(\frac{t}{\unsure{p\rho|p^\rho}}\right)
\]
% archive sheet number: 8

